% 02-resultados-pt.tex — Resultados
\section{Resultados}
\label{sec:results}

\subsection{Ranking de problema único (R499)}
\label{sec:r499}
Sob C2 no problema-prova NT-001, os resultados do primeiro trial foram:
mimo 0,987 (23,8\,s), big-pickle 0,923 (43,0\,s), nemotron 0,709 (107,4\,s),
muse-1,3 0,324, muse-1,2 0,312. Os modelos muse-spark nunca entregaram
resposta válida em nenhuma condição (``raciocínio falso''). Saídas vazias
observadas em C0 e C1; nemotron timeout de 110\,s em C1.

\subsection{Validação cruzada após reinício do host (R503)}
\label{sec:r503}
A Tabela~\ref{tab:acc} e as Figuras~\ref{fig:acc}--\ref{fig:lat} apresentam
o benchmark 9 $\times$ 3 executado após o reinício do host.

\begin{table}[htbp]
\centering
\caption{Acurácia validada cruzada ($n=9$), IC Clopper--Pearson 95\%,
timeouts, latências (chamadas ok).}
\label{tab:acc}
\begin{tabular}{lccccc}
\toprule
Modelo & acertos & acurácia & IC 95\% & timeouts & latência (s) \\
\midrule
mimo-v2.5-free    & 2/9 & 0,22 & [0,03; 0,60] & 2 & 43,5 (29,4--91,3) \\
big-pickle        & 1/9 & 0,11 & [0,00; 0,48] & 0 & 31,3 (20,0--49,4) \\
nemotron-3-ultra  & 2/9 & 0,22 & [0,03; 0,60] & 0 & 92,5 (75,0--136,7) \\
\bottomrule
\end{tabular}
\end{table}

\begin{figure}[htbp]
\centering
\includegraphics[width=0.85\linewidth]{figures/fig1_accuracies.pdf}
\caption{Acurácia com IC Clopper--Pearson 95\% (barras de erro), $n=9$.}
\label{fig:acc}
\end{figure}

\begin{figure}[htbp]
\centering
\includegraphics[width=0.7\linewidth]{figures/fig2_matrix.pdf}
\caption{Matriz de acertos: 9 problemas $\times$ 3 modelos (check = correto).}
\label{fig:matrix}
\end{figure}

\begin{figure}[htbp]
\centering
\includegraphics[width=0.85\linewidth]{figures/fig3_latency.pdf}
\caption{Distribuição de latência por modelo; losango vermelho = timeout.}
\label{fig:lat}
\end{figure}

\subsection{Estatísticas pareadas}
\label{sec:paired}
Testes McNemar com correção de continuidade são não-significativos com $n=9$
(mimo vs.\ big-pickle: $b=1, c=0$, $\chi^2=0,00$, $p=1,00$; mimo vs.\
nemotron: $b=1, c=1$, $\chi^2=0,50$, $p=0,48$; big-pickle vs.\ nemotron:
$b=0, c=1$, $\chi^2=0,00$, $p=1,00$). Com $n=9$ o poder estatístico é
baixo; reportamos $p$-valores exatos. O $\kappa$ de Cohen entre mimo e
big-pickle é $0,609$ (concordância substancial --- ambos acertam o mesmo
problema NT-002), enquanto entre mimo e nemotron é $0,357$ (concordância
razoável).

\subsection{Variabilidade entre execuções}
\label{sec:variabilidade}
Comparando duas execuções independentes do mesmo protocolo (mesmo host,
mesmos parâmetros):
\begin{itemize}
  \item \textbf{Execução 1} (pós-reinício imediato): mimo 4/9 (0,44),
  big-pickle 1/9 (0,11), nemotron 1/9 (0,11).
  \item \textbf{Execução 2} (horas depois): mimo 2/9 (0,22), big-pickle 1/9
  (0,11), nemotron 2/9 (0,22).
\end{itemize}
A variabilidade é substancial (mimo variou de 0,44 para 0,22 entre
execuções). Isso reforça a necessidade de múltiplos trials por célula em
benchmarks de LLMs --- uma prática padrão em outros campos mas incomum em
relatórios de modelos de linguagem.

\subsection{Achados relevantes para metodologia}
\label{sec:achados}
\textbf{(a) Vazamento no arsenal original} (R442): o resolvedor padrão ecoava
a resposta e auto-pontuava 7. O resolvedor determinístico (R499) reduziu
para escores reais (2/4 canônicos, média 4,5/7; corpus aumentado 8/9 com
``parcial'' explícito para algebra-004). \textbf{(b) Sensibilidade de
serialização}: o \emph{grading head} é sensível a $2^{(n-1)}$ vs.\
$2^{n-1}$ (combinatoria-001). \textbf{(c) Hipótese do estado do host
refutada}: a replicação R500 atribuía as saídas vazias do big-pickle ao
estado acumulado do host; após reinício completo, o big-pickle ainda falha
em 8/9 problemas. \textbf{(d) Roteamento (R500)}: \texttt{route\_free
("academic")} retorna mimo/0,987 com fallbacks explícitos; o catálogo free
agora inclui \texttt{big-pickle}; 46 testes verdes.

\subsection{Regressão e gates}
A suíte completa de testes passa: 4077 aprovados, 58 ignorados, 0 falhas
(577\,s, pré-R503), mais 11 testes R503+R500 aprovados (1,7\,s), e 26
testes R504+R505 aprovados (0,6\,s).

\subsection{Benchmark de código executável (R506)}
\label{sec:r506}
Para avaliar se os modelos realmente \emph{resolvem} os problemas (e não
apenas geram texto que parece resposta), implementamos o pipeline R506:
cada LLM recebe um prompt que solicita um script Python completo, o código
é extraído, executado em sandbox seguro (subprocess com timeout de 30\,s,
sem acesso a rede), e a saída \texttt{stdout} é comparada com a resposta
esperada. A Tabela~\ref{tab:r506} apresenta os resultados.

\begin{table}[htbp]
\centering
\caption{R506 --- Benchmark de código executável ($n=9$). Status:
\textbf{C} = correto, \textbf{W} = resposta errada, \textbf{E} = erro
de runtime, \textbf{M} = código ausente.}
\label{tab:r506}
\begin{tabular}{lccccc}
\toprule
Modelo & C & W & E & M & acurácia \\
\midrule
mimo-v2.5-free    & 1 & 1 & 2 & 5 & 0,111 \\
big-pickle        & 0 & 1 & 2 & 6 & 0,000 \\
nemotron-3-ultra  & 2 & 1 & 4 & 2 & 0,222 \\
\bottomrule
\end{tabular}
\end{table}

\begin{figure}[htbp]
\centering
\includegraphics[width=0.85\linewidth]{figures/fig4_code_vs_text.pdf}
\caption{Comparação R503 (texto) vs R506 (código executável). A capacidade
de gerar código correto é consistentemente inferior à de gerar texto que
parece resposta.}
\label{fig:code_vs_text}
\end{figure}

Os resultados revelam três achados críticos. \emph{Primeiro}, a taxa de
``código ausente'' (modelos que não geram nenhum bloco de código) é alta:
5/9 para mimo, 6/9 para big-pickle e 2/9 para nemotron. Isso indica que
muitos modelos respondem em linguagem natural quando solicitados a gerar
código, ou geram blocos de formato incorreto. \emph{Segundo}, entre os
códigos gerados, a taxa de erro de runtime é substancial: 2/4 para mimo,
2/3 para big-pickle e 4/7 para nemotron --- os modelos geram código que
\emph{não executa}. \emph{Terceiro}, a comparação R503 vs R506 é
iluminadora: mimo cai de 0,222 (texto) para 0,111 (código); big-pickle
cai de 0,111 para 0,000; nemotron permanece em 0,222. A capacidade de
gerar código executável correto é \emph{consistente e mensuravelmente
inferior} à de gerar texto que parece conter a resposta.

O nemotron se destaca por gerar mais código (7/9 problemas, média 317\,B)
mas com alta taxa de erro de runtime (4/7), indicando ``confiança
descalibrada'' --- o modelo gera código elaborado que falha na execução.
O big-pickle gera pouco código (3/9) mas quando gera, também falha (2/3
runtime errors). O mimo gera código intermediário (4/9) com taxa de erro
moderada (2/4).